\begin{theorem}[Disjunction property]
\label{thm:disjunction-property}
If \(\LJ\derives\;\Rightarrow A\lor B\), then
\(\LJ\derives\;\Rightarrow A\) or \(\LJ\derives\;\Rightarrow B\).
\textit{Status: complete proof.}
\end{theorem}

\begin{proof}
By cut elimination for \(\LJ\) (\cref{ch:cut-elimination}), fix a
cut-free \(\LJ\)-derivation of \(\Rightarrow A\lor B\).  Its final
inference is not \textsc{Id} or \(\bottom L\), whose conclusions have
nonempty antecedents, and not a left logical or left structural rule, for
the same reason.  The calculus has no right weakening or right contraction.
Among the remaining right rules, only \(\lor R_1\) and \(\lor R_2\)
conclude a sequent whose formula is a disjunction.  Hence the final
inference has premise \(\Rightarrow A\) or \(\Rightarrow B\).
\end{proof}

Classical derivability lacks this property: \(P\lor\neg P\) is
classically derivable while neither disjunct is.

